
\documentclass[12pt, a4paper]{amsart}


\usepackage{amsmath,amssymb,amscd,amsfonts}
\begin{document}
\centerline{\bf MATHEMATICS 4400}
\centerline{\bf Math 4400 Intro To Number Theory}
\centerline{Fall semester 2022}
\bigskip

  {\bf Text:}{  Numbers, Groups and Cryptography, by G. Savin}
 
{\bf When/where:}{ M, W, F 10.45-11.35  A.M. LCB 225}
  
{\bf instructor:}{ Christopher Hacon}
  
{\bf  office:} JWB 301

 {\bf telephone:} 801-581-7429
  
{\bf email:}{hacon(at)math.utah.edu}

 {\bf office hours:} M 9:45-10:35, W 9:45-10:35, F 12:40-1:20.

\bigskip
\noindent
{\bf Course outline:}
 	
An overview of algebraic number theory, covering factorization and primes, modular arithmetic, quadratic residues, continued fractions, quadratic forms, and diophantine equations

\noindent
{\bf Grading:}\qquad There will be weekly quizzes, and a comprehensive final 
examination.  Quizzes will account for 60\% of your grade 
(you will be 
allowed to drop your 3 lowest scores, however there will be no make up tests),  
and the final exam will make up the 
remaining 40\%.


Labor day Holiday Sept 5; Fall break Sun-Sun, October 9-16; Thanksgiving break 
Thurs.-Fri., Nov. 24-27

\noindent
{\bf Final:} Thursday, Dec. 15, 2022
10:30 am - 12:30 pm

\medskip\noindent
{\bf Homepage:} You may refer to the course homepage
for further information and future handouts.

\noindent
{ {http://www.math.utah.edu/\~ {} hacon/4400.html}}


\medskip\noindent
{\bf ADA statement:}  The American with Disabilities Act requires 
that reasonable accommodations be provided for students with physical, 
sensory, 
cognitive, systemic, learning, and psychiatric disabilities.  
Please contact me at the beginning of the semester to discuss 
any such accommodations for the course.


\end{document}
\end
























































































